\chapter{Perfect Completeness of the IPA Scheme}

This chapter states and sketches the first of the two headline security
properties of the IPA polynomial commitment of \cref{def:ipa_scheme}:
\emph{perfect completeness}. Informally, an honest prover always convinces the
verifier of a true evaluation: committing to a polynomial \(p\) of degree
\(< d\) and then running the honest opening protocol on a point \(x\) with the
true value \(v = p(x)\) makes the verifier accept with probability \(1\). This
is the completeness half of Halo's Theorem~1; binding and extractability are
deferred to their own chapter.

The Lean statement follows the two-layer shape used by ArkLib's worked KZG
development. There, an \emph{algebraic} correctness lemma
(\texttt{KZG.correctness}), asserting that the concrete commit/open/verify
operations satisfy the verification equation, is \emph{lifted} to the
interface-level predicate \texttt{Commitment.perfectCorrectness}
(\texttt{KZG.CommitmentScheme.correctness}) by unfolding the oracle-computation
plumbing of \texttt{KeyGen}, \texttt{Commit}, and the opening reduction and
discharging the probability-one obligation. The IPA correctness theorem below
plays the role of that interface-level statement: its proof reduces, after the
same unfolding, to the algebraic accept-equation identity sketched here for the
honest transcript of \cref{def:ipa_opening_protocol}.

\section{Perfect completeness}

\begin{theorem}[Perfect completeness of IPA]
  \label{thm:ipa_correctness}
  \lean{Kimchi.Commitment.IPA.Correctness.ipa_correctness}
  \uses{def:ipa_scheme, def:commitment_perfect_correctness,
    def:ipa_opening_protocol, def:ipa_commit}
  % SOURCE: [Halo], Section 3, Theorem 1 (read from references/halo.txt)
  % SOURCE QUOTE: "Theorem 1. The protocol presented in Section 3.1 has perfect completeness,
  % computational witness-extended emulation, and perfect special honest-verifier
  % zero knowledge."
  \textit{Source: Halo, Section 3, Theorem 1 (perfect completeness).}
  The IPA commitment scheme \(\texttt{ipa}\) of \cref{def:ipa_scheme} satisfies
  perfect correctness in the sense of \cref{def:commitment_perfect_correctness}:
  for every polynomial \(p \in \Fr[X]\) of degree \(< d\), presented by its
  coefficient vector \(a \in \Fr^{d}\), and every evaluation point
  \(x \in \Fr\), the experiment that
  \begin{enumerate}
    \item runs key generation to obtain the SRS \(\sigma = (g, H)\),
    \item runs the honest \(\Commit\) of \cref{def:ipa_commit} on \(a\),
      yielding the commitment \(P = \inner{a}{g} + [r]H\) with decommitment
      \(r\), and
    \item runs the honest opening protocol of
      \cref{def:ipa_opening_protocol} on the statement
      \((P, x, v)\) with \(v = p(x)\),
  \end{enumerate}
  makes the verifier accept with probability \(1\); equivalently,
  \(\texttt{ipa}\) satisfies \(\correctness\) with error
  \(\varepsilon_{\mathrm{c}} = 0\).
\end{theorem}

% SOURCE QUOTE PROOF: "and the prover proceeds
%                     Pk to demonstrate   Pk knowledge of a = a 0 and synthetic
% blinding factor r0 = j=1 (lj u2j ) + r + j=1 (rj u−2
%                                                   j  ) such that
%                              Q = [a] G + [r0 ]H + [ab] U
%                                                                                 (2)
%                                   = [a] (G + [b] U ) + [r0 ]H
% which establishes the claim that v = v 0 = ha, bi as described in [13, Section 3]."
% SOURCE QUOTE PROOF: "The prover begins by sampling random d, s ∈ Fp and sending
%
%                               R = [d] (G + [b] U ) + [s]H
%
% which the verifier responds to with random challenge c ∈ I. The prover now
% sends
%                                  z1 = ac + d
%                                  z2 = cr0 + s
% and the verifier accepts if [c] Q + R = [z1 ] (G + [b] U ) + [z2 ]H."
\begin{proof}
  \uses{def:ipa_opening_proof, def:ipa_bpoly, def:ipa_relation,
    def:mathlib_module_smul}
  \textit{Source: Halo, Section 3.1 and eq.~(2).}
  The proof mirrors the perfect-completeness half of Halo's Theorem~1: it tracks
  the honest prover's values through the protocol of
  \cref{def:ipa_opening_protocol} and verifies the final accept equation by
  direct substitution. Since the honest \(\Commit\) makes the statement
  \((P, x, v)\) a member of the opening relation of \cref{def:ipa_relation}
  (i.e.\ \(P = \inner{a}{g} + [r]H\) and \(v = p(x) = \inner{a}{(1, x, \dots,
  x^{d-1})}\)), every step below is an identity in \(\GG\) and \(\Fr\), so the
  accepting event has probability \(1\). At the interface level this algebraic
  identity is what the \texttt{perfectCorrectness} unfolding reduces to, exactly
  as ArkLib's \texttt{KZG.CommitmentScheme.correctness} reduces to its algebraic
  \texttt{KZG.correctness} lemma.

  \emph{(i) Collapse of the fold.} After the setup move both parties hold
  \(P' = P + [v]U\), and the prover initializes \(g' := g\), \(a' := a\),
  \(b' := (1, x, x^{2}, \dots, x^{d-1})\). Each of the \(k = \log_{2} d\) honest
  folding rounds replaces a length-\(2m\) triple \((g', a', b')\) by the
  length-\(m\) triple
  \[
    a' \leftarrow a'_{\mathrm{hi}}\, u_{j}^{-1} + a'_{\mathrm{lo}}\, u_{j},
    \quad
    b' \leftarrow b'_{\mathrm{lo}}\, u_{j}^{-1} + b'_{\mathrm{hi}}\, u_{j},
    \quad
    g' \leftarrow g'_{\mathrm{lo}}\, u_{j}^{-1} + g'_{\mathrm{hi}}\, u_{j}.
  \]
  After all \(k\) rounds the three vectors have length one. Write
  \(a_{0}\) for the final folded scalar \(a'_{0}\), \(g_{0} := g'_{0}\), and
  \(b_{0} := b'_{0}\). Unwinding the recursion expresses \(g_{0}\) and \(b_{0}\)
  as the inner products of the length-\(d\) coefficient vector \(s\) of
  \cref{def:ipa_bpoly} (eq.~(1)) against \(g\) and against the powers of \(x\):
  \[
    g_{0} = \inner{s}{g} =: G,
    \qquad
    b_{0} = \inner{s}{(1, x, x^{2}, \dots, x^{d-1})} = \bpoly(x) =: b,
  \]
  the second equality being the defining property of the challenge polynomial
  \(\bpoly\) of \cref{def:ipa_bpoly}.

  \emph{(ii) Telescoping to \(Q\) (eq.~(2)).} The verifier forms
  \[
    Q = \sum_{j=1}^{k} [u_{j}^{2}]\, L_{j} \;+\; P' \;+\;
        \sum_{j=1}^{k} [u_{j}^{-2}]\, R_{j}.
  \]
  Substituting the honest cross-term definitions
  \(L_{j} = \inner{a'_{\mathrm{lo}}}{g'_{\mathrm{hi}}} + [l_{j}]H +
  [\inner{a'_{\mathrm{lo}}}{b'_{\mathrm{hi}}}]U\) and
  \(R_{j} = \inner{a'_{\mathrm{hi}}}{g'_{\mathrm{lo}}} + [r_{j}]H +
  [\inner{a'_{\mathrm{hi}}}{b'_{\mathrm{lo}}}]U\) of round \(j\), the
  cross-terms telescope against \(P'\): each round's contribution combines with
  the running \(\inner{a'}{g'}\)- and \(\inner{a'}{b'}\)-terms so that all
  intermediate pieces cancel, leaving the length-one residue. With the synthetic
  blinding factor
  \(r' = \sum_{j=1}^{k} (l_{j} u_{j}^{2}) + r + \sum_{j=1}^{k} (r_{j}
  u_{j}^{-2})\) one obtains Halo's eq.~(2):
  \[
    Q = [a_{0}]\,G + [r']H + [a_{0} b]\,U
      = [a_{0}]\,(G + [b]U) + [r']H.
  \]

  \emph{(iii) The Schnorr check.} The prover samples \(d, s \in \Fr\), sends
  \(\delta = [d](G + [b]U) + [s]H\) (Halo's \(R\)), receives the Schnorr
  challenge \(c \in \Fr\), and answers with the honest responses
  \(z_{1} = a_{0} c + d\) and \(z_{2} = c r' + s\). Substituting these and the
  expressions for \(Q\) and \(\delta\) into the right-hand side of the
  verifier's check and using bilinearity of the scalar action
  (\cref{def:mathlib_module_smul}),
  \begin{align*}
    [z_{1}]\,(G + [b]U) + [z_{2}]\,H
      &= [a_{0} c + d]\,(G + [b]U) + [c r' + s]\,H \\
      &= [c]\bigl([a_{0}]\,(G + [b]U) + [r']H\bigr)
         + \bigl([d]\,(G + [b]U) + [s]H\bigr) \\
      &= [c]\,Q + \delta.
  \end{align*}
  This is exactly the verifier's accept condition
  \([c]\,Q + \delta = [z_{1}]\,(G + [b]U) + [z_{2}]\,H\), so the honest
  transcript \((\{(L_{j}, R_{j})\}_{j}, a_{0}, g_{0}, \delta, z_{1}, z_{2})\) of
  \cref{def:ipa_opening_proof} always passes. Hence the verifier accepts with
  probability \(1\), which is perfect completeness.
\end{proof}
